\documentclass{article}
\usepackage{amsmath,amssymb,amsthm}
\usepackage{algorithm}
\usepackage[noend]{algorithmic}
\usepackage{enumitem}
\usepackage{hyperref}
\usepackage{bm}

\newtheorem{theorem}{Theorem}[section]
\newtheorem{lemma}{Lemma}[section]
\newtheorem{assumption}{Assumption}[section]

\newcommand{\EE}{\mathbb{E}}
\newcommand{\RR}{\mathbb{R}}
\newcommand{\norm}[1]{\left\|#1\right\|}
\newcommand{\inner}[2]{\left\langle #1,#2\right\rangle}
\newcommand{\grad}{\nabla}
\newcommand{\prox}{\operatorname{prox}}

\begin{document}

\section*{Algorithm Name}
\textbf{Stochastic Variance Reduced Gradient (SVRG) for Non‑Convex Smooth Optimization}

\section*{Mathematical Formulation}
We consider a finite‑sum objective
\[
f(\mathbf{x})\;=\;\frac{1}{n}\sum_{i=1}^{n} f_i(\mathbf{x}),\qquad 
\mathbf{x}\in\RR^{d},
\]
where each component $f_i$ is differentiable and $L$–smooth:
\[
\forall \mathbf{u},\mathbf{v}\in\RR^{d},\qquad 
f_i(\mathbf{u})\le f_i(\mathbf{v})+\inner{\grad f_i(\mathbf{v})}{\mathbf{u}-\mathbf{v}}+
\frac{L}{2}\norm{\mathbf{u}-\mathbf{v}}^{2}. \tag{1}
\]

SVRG proceeds in \emph{epochs}.  At the beginning of epoch $s$ a \emph{snapshot}
$\tilde{\mathbf{x}}^{\,s}$ is fixed and its full gradient is computed:
\[
\tilde{\bm\mu}^{\,s}\;:=\;\grad f(\tilde{\mathbf{x}}^{\,s})=
\frac{1}{n}\sum_{i=1}^{n}\grad f_i(\tilde{\mathbf{x}}^{\,s}). \tag{2}
\]

Inside epoch $s$ we perform $m$ inner stochastic updates.  
For inner iteration $t=0,\dots,m-1$ we draw an index $i_t$ uniformly
from $\{1,\dots,n\}$ and construct the \emph{variance‑reduced estimator}
\[
\mathbf{v}_{t}
\;:=\;\grad f_{i_t}(\mathbf{x}_{t})-
\grad f_{i_t}(\tilde{\mathbf{x}}^{\,s})+\tilde{\bm\mu}^{\,s}. \tag{3}
\]
The iterate is then updated by a standard gradient step
\[
\mathbf{x}_{t+1}\;=\;\mathbf{x}_{t}-\eta\,\mathbf{v}_{t}, \qquad
\eta>0\text{ (stepsize).} \tag{4}
\]

After $m$ inner steps we set the next snapshot
$\tilde{\mathbf{x}}^{\,s+1}\gets\mathbf{x}_{m}$ (or any convex combination of the
inner iterates) and start a new epoch.

\section*{Pseudocode}
\begin{algorithm}[H]
\caption{SVRG for non‑convex $L$‑smooth finite‑sum problems}
\label{alg:svrg}
\begin{algorithmic}[1]
\REQUIRE Initial point $\tilde{\mathbf{x}}^{\,0}$, stepsize $\eta$, inner loop length $m$
\FOR{epoch $s=0,1,2,\dots$}
    \STATE Compute full gradient $\tilde{\bm\mu}^{\,s}\gets\grad f(\tilde{\mathbf{x}}^{\,s})$ \label{alg:fullgrad}
    \STATE Set $\mathbf{x}_{0}\gets\tilde{\mathbf{x}}^{\,s}$
    \FOR{inner iteration $t=0$ \TO $m-1$}
        \STATE Sample $i_{t}\sim\text{Uniform}\{1,\dots,n\}$
        \STATE $\mathbf{v}_{t}\gets\grad f_{i_t}(\mathbf{x}_{t})-
               \grad f_{i_t}(\tilde{\mathbf{x}}^{\,s})+\tilde{\bm\mu}^{\,s}$ \label{alg:vr}
        \STATE $\mathbf{x}_{t+1}\gets\mathbf{x}_{t}-\eta\,\mathbf{v}_{t}$ \label{alg:update}
    \ENDFOR
    \STATE Set $\tilde{\mathbf{x}}^{\,s+1}\gets\mathbf{x}_{m}$  \COMMENT{or any average}
\ENDFOR
\end{algorithmic}
\end{algorithm}

\section*{Dry Run Example}
Consider a single‑sample problem ($n=1$) with
\[
f(w)=(w-3)^{2},\qquad \grad f(w)=2(w-3).
\]
Take $w_{0}=0$, stepsize $\eta=0.1$, inner‑loop length $m=3$.
Since $n=1$, the full gradient equals the stochastic gradient at every point,
hence the variance‑reduced estimator (3) reduces to the ordinary gradient.

\begin{center}
\begin{tabular}{c|c|c|c}
$t$ & $w_{t}$ & $\grad f(w_{t})$ & $f(w_{t})$ \\ \hline
0 & $0$ & $2(0-3)=-6$ & $9$ \\ 
1 & $w_{1}=0-0.1(-6)=0.6$ & $2(0.6-3)=-4.8$ & $(0.6-3)^{2}=5.76$ \\ 
2 & $w_{2}=0.6-0.1(-4.8)=1.08$ & $2(1.08-3)=-3.84$ & $(1.08-3)^{2}=3.6864$ \\ 
3 & $w_{3}=1.08-0.1(-3.84)=1.464$ & $2(1.464-3)=-3.072$ & $(1.464-3)^{2}=2.366$ 
\end{tabular}
\end{center}
The objective value decreased from $9$ to $\approx2.37$ after three SVRG inner steps,
illustrating the same behaviour as ordinary SGD in this trivial case.

\newpage
\section*{Assumptions}
\begin{assumption}[Smoothness]\label{ass:smooth}
Each component $f_i$ is $L$‑smooth, i.e.\ inequality (1) holds for a common $L>0$.
\end{assumption}
\begin{assumption}[Finite Sum]\label{ass:finite}
The objective is a finite average $f(\mathbf{x})=\frac1n\sum_{i=1}^{n}f_i(\mathbf{x})$.
\end{assumption}
\begin{assumption}[Bounded Variance of Stochastic Gradients]\label{ass:var}
For all $\mathbf{x}$,
\[
\EE_{i\sim\mathcal{U}[n]}\!\big[\norm{\grad f_i(\mathbf{x})-\grad f(\mathbf{x})}^{2}\big]
\;\le\;\sigma^{2},
\]
with a constant $\sigma\ge0$.
\end{assumption}
\begin{assumption}[Stepsize]\label{ass:stepsize}
The stepsize satisfies $0<\eta\le\frac{1}{3L}$.
\end{assumption}
All expectations $\EE[\cdot]$ are taken over the random choice of indices
$\{i_t\}$ conditioned on the current outer snapshot.

\section*{Main Convergence Theorem}
\begin{theorem}\label{thm:main}
Assume \ref{ass:smooth}–\ref{ass:stepsize} hold.  
Let $\{\tilde{\mathbf{x}}^{\,s}\}_{s\ge0}$ be the sequence of snapshots generated by
Algorithm~\ref{alg:svrg} with inner loop length $m\ge1$.
Define the total number of stochastic gradient evaluations
$N:=S\cdot(m+n)$ after $S$ epochs.
Then for any $S\ge1$,
\[
\frac{1}{S\,m}\sum_{s=0}^{S-1}\sum_{t=0}^{m-1}
\EE\!\big[\norm{\grad f(\mathbf{x}_{t}^{\,s})}^{2}\big]
\;\le\;
\frac{2\big(f(\tilde{\mathbf{x}}^{\,0})-f^{\star}\big)}{\eta\,S\,m}
\;+\;4L\eta\,\sigma^{2},
\tag{5}
\]
where $f^{\star}:=\inf_{\mathbf{x}}f(\mathbf{x})$.
Consequently, choosing $\eta= \Theta\!\big(\frac{1}{L\sqrt{S}}\big)$ yields
\[
\EE\!\big[\norm{\grad f(\mathbf{x}_{\mathrm{out}})}^{2}\big]
\;=\;O\!\Big(\frac{L\big(f(\tilde{\mathbf{x}}^{\,0})-f^{\star}\big)}{N}\Big)
\;+\;O\!\Big(\frac{\sigma^{2}}{\sqrt{N}}\Big),
\]
i.e. the expected squared gradient norm decays as $O(1/N)$ up to the
variance term.
\end{theorem}

\section*{Theoretical Properties}
\begin{itemize}[leftmargin=*]
\item \textbf{Stability.} The variance‑reduced estimator $\mathbf{v}_t$ is unbiased
      (Lemma~\ref{lem:unbiased}) and its second moment is uniformly bounded
      (Lemma~\ref{lem:variance}), guaranteeing that the iterates do not
      diverge provided $\eta\le 1/(3L)$.
\item \textbf{Hyperparameter Sensitivity.}
      The bound (5) shows a linear dependence on $\eta$ for the variance term
      and a $1/\eta$ dependence for the bias term.
      The optimal trade‑off is $\eta\asymp 1/(L\sqrt{S})$, which matches the
      classical SGD rate for non‑convex smooth problems.
\item \textbf{Comparison to Baselines.}
      Ordinary SGD with stepsize $\eta=O(1/\sqrt{N})$ achieves
      $O(1/\sqrt{N})$ in the non‑convex setting, whereas SVRG attains the
      faster $O(1/N)$ rate (up to the additive variance term) because the
      control variate $\tilde{\bm\mu}^{\,s}$ reduces stochastic noise.
\end{itemize}

\section*{Discussion}
The theorem demonstrates that even without convexity, the double‑loop
SVRG scheme enjoys a \emph{linear} convergence in the sense of
expected gradient norm decay, provided the stepsize is sufficiently small.
The result hinges on the $L$‑smoothness of each component and on the
unbiasedness of the variance‑reduced estimator.
When $\sigma=0$ (i.e.\ each $f_i$ already shares the same gradient), the
method reduces to full‑gradient descent and the bound simplifies to the
deterministic $O(1/T)$ rate.  In practice, modest inner loop lengths
($m\approx 2n$) are sufficient to reap the theoretical gains.

\newpage
\section*{Preliminaries (Definitions and Lemmas)}
\begin{lemma}[Unbiasedness]\label{lem:unbiased}
For any inner iteration $t$ and any outer snapshot $\tilde{\mathbf{x}}$,
\[
\EE_{i_t}\!\big[\mathbf{v}_t\mid \mathbf{x}_t,\tilde{\mathbf{x}}\big]
\;=\;\grad f(\mathbf{x}_t).
\]
\end{lemma}
\begin{proof}
\[
\EE_{i_t}\!\big[\mathbf{v}_t\mid\cdot\big]
=\frac{1}{n}\sum_{i=1}^{n}\big(\grad f_i(\mathbf{x}_t)-\grad f_i(\tilde{\mathbf{x}})\big)
+\grad f(\tilde{\mathbf{x}})
=\grad f(\mathbf{x}_t). \quad\blacksquare
\]
\end{proof}

\begin{lemma}[Second‑moment bound]\label{lem:variance}
Under Assumptions~\ref{ass:smooth} and \ref{ass:var},
\[
\EE\!\big[\norm{\mathbf{v}_t}^{2}\mid \mathbf{x}_t,\tilde{\mathbf{x}}\big]
\;\le\;
4L\big(f(\mathbf{x}_t)-f^{\star}\big)
\;+\;
4L\big(f(\tilde{\mathbf{x}})-f^{\star}\big)
\;+\;
2\sigma^{2}.
\tag{6}
\]
\end{lemma}
\begin{proof}
Start from
\[
\mathbf{v}_t
=\grad f_{i_t}(\mathbf{x}_t)-\grad f_{i_t}(\tilde{\mathbf{x}})+\tilde{\bm\mu},
\qquad
\tilde{\bm\mu}=\grad f(\tilde{\mathbf{x}}).
\]
Add and subtract $\grad f(\mathbf{x}_t)$:
\[
\mathbf{v}_t
=\big(\grad f_{i_t}(\mathbf{x}_t)-\grad f(\mathbf{x}_t)\big)
-\big(\grad f_{i_t}(\tilde{\mathbf{x}})-\grad f(\tilde{\mathbf{x}})\big)
+\grad f(\mathbf{x}_t).
\]
Take squared norm and use $(a+b)^{2}\le 2a^{2}+2b^{2}$:
\[
\norm{\mathbf{v}_t}^{2}
\le 2\norm{\grad f(\mathbf{x}_t)}^{2}
   +2\norm{\big(\grad f_{i_t}(\mathbf{x}_t)-\grad f(\mathbf{x}_t)\big)
          -\big(\grad f_{i_t}(\tilde{\mathbf{x}})-\grad f(\tilde{\mathbf{x}})\big)}^{2}.
\tag{7}
\]
Apply the inequality $\norm{a-b}^{2}\le2\norm{a}^{2}+2\norm{b}^{2}$ to the second term:
\[
\begin{aligned}
&\norm{\big(\grad f_{i_t}(\mathbf{x}_t)-\grad f(\mathbf{x}_t)\big)
          -\big(\grad f_{i_t}(\tilde{\mathbf{x}})-\grad f(\tilde{\mathbf{x}})\big)}^{2}\\
&\le2\norm{\grad f_{i_t}(\mathbf{x}_t)-\grad f(\mathbf{x}_t)}^{2}
        +2\norm{\grad f_{i_t}(\tilde{\mathbf{x}})-\grad f(\tilde{\mathbf{x}})}^{2}.
\end{aligned}
\tag{8}
\]
Take expectation w.r.t.\ $i_t$ and use Assumption~\ref{ass:var}:
\[
\EE_{i_t}\!\big[\norm{\grad f_{i_t}(\mathbf{x}_t)-\grad f(\mathbf{x}_t)}^{2}\big]
\le\sigma^{2},\qquad
\EE_{i_t}\!\big[\norm{\grad f_{i_t}(\tilde{\mathbf{x}})-\grad f(\tilde{\mathbf{x}})}^{2}\big]
\le\sigma^{2}.
\tag{9}
\]
Now bound $\norm{\grad f(\mathbf{x}_t)}^{2}$ using $L$‑smoothness.
For any $L$‑smooth function,
\[
\frac{1}{2L}\norm{\grad f(\mathbf{x})}^{2}
\;\le\;f(\mathbf{x})-f^{\star},
\tag{10}
\]
which follows from integrating the gradient inequality.
Thus,
\[
\norm{\grad f(\mathbf{x}_t)}^{2}\le2L\big(f(\mathbf{x}_t)-f^{\star}\big),
\qquad
\norm{\grad f(\tilde{\mathbf{x}})}^{2}\le2L\big(f(\tilde{\mathbf{x}})-f^{\star}\big).
\tag{11}
\]
Combining (7)–(11) and multiplying by the outer expectation yields (6). \quad\blacksquare
\end{proof}

\section*{Main Proof (Step‑by‑Step Derivation)}
We prove Theorem~\ref{thm:main}.  Throughout the proof we fix an epoch $s$
and suppress the superscript $s$ for readability (i.e.\ $\tilde{\mathbf{x}}:=\tilde{\mathbf{x}}^{\,s}$,
$\tilde{\bm\mu}:=\tilde{\bm\mu}^{\,s}$).

\noindent\textbf{[Step 1] Smoothness descent inequality.}
From $L$‑smoothness (1) applied to $f$,
\[
f(\mathbf{x}_{t+1})
\le f(\mathbf{x}_{t})+\inner{\grad f(\mathbf{x}_{t})}{\mathbf{x}_{t+1}-\mathbf{x}_{t}}
+\frac{L}{2}\norm{\mathbf{x}_{t+1}-\mathbf{x}_{t}}^{2}.
\tag{12}
\]
Using the update rule (4), $\mathbf{x}_{t+1}-\mathbf{x}_{t}=-\eta\mathbf{v}_{t}$,
so (12) becomes
\[
f(\mathbf{x}_{t+1})
\le f(\mathbf{x}_{t})-\eta\inner{\grad f(\mathbf{x}_{t})}{\mathbf{v}_{t}}
+\frac{L\eta^{2}}{2}\norm{\mathbf{v}_{t}}^{2}.
\tag{13}
\]

\noindent\textbf{[Step 2] Take conditional expectation.}
Condition on $\mathbf{x}_{t}$ and $\tilde{\mathbf{x}}$ and apply Lemma~\ref{lem:unbiased}:
\[
\EE\!\big[\inner{\grad f(\mathbf{x}_{t})}{\mathbf{v}_{t}}\mid\mathbf{x}_{t},\tilde{\mathbf{x}}\big]
=\norm{\grad f(\mathbf{x}_{t})}^{2}.
\tag{14}
\]
Similarly, keep the second term as is for now:
\[
\EE\!\big[\norm{\mathbf{v}_{t}}^{2}\mid\mathbf{x}_{t},\tilde{\mathbf{x}}\big]
\le 4L\big(f(\mathbf{x}_{t})-f^{\star}\big)
   +4L\big(f(\tilde{\mathbf{x}})-f^{\star}\big)
   +2\sigma^{2}
\tag{15}
\]
by Lemma~\ref{lem:variance}.

Insert (14)–(15) into (13) and then take full expectation:
\[
\EE\!\big[f(\mathbf{x}_{t+1})\big]
\le \EE\!\big[f(\mathbf{x}_{t})\big]
-\eta\EE\!\big[\norm{\grad f(\mathbf{x}_{t})}^{2}\big]
+\frac{L\eta^{2}}{2}\EE\!\big[\,4L\!\big(f(\mathbf{x}_{t})-f^{\star}\big)
   +4L\!\big(f(\tilde{\mathbf{x}})-f^{\star}\big)+2\sigma^{2}\big].
\tag{16}
\]

\noindent\textbf{[Step 3] Rearrange terms.}
Define $A_{t}:=\EE\!\big[f(\mathbf{x}_{t})-f^{\star}\big]$ and note that
$f(\tilde{\mathbf{x}})-f^{\star}$ is constant throughout the inner loop.
Equation (16) becomes
\[
A_{t+1}
\le A_{t}
-\eta\EE\!\big[\norm{\grad f(\mathbf{x}_{t})}^{2}\big]
+2L^{2}\eta^{2}A_{t}
+2L^{2}\eta^{2}\big(f(\tilde{\mathbf{x}})-f^{\star}\big)
+L\eta^{2}\sigma^{2}.
\tag{17}
\]

\noindent\textbf{[Step 4] Choose stepsize $\eta\le\frac{1}{3L}$.}
Since $2L^{2}\eta^{2}\le\frac{2}{9}$, we have $1-2L^{2}\eta^{2}\ge\frac{7}{9}>0$.
Move the $2L^{2}\eta^{2}A_{t}$ term to the left side:
\[
(1-2L^{2}\eta^{2})A_{t+1}
\le A_{t}
-\eta\EE\!\big[\norm{\grad f(\mathbf{x}_{t})}^{2}\big]
+2L^{2}\eta^{2}\big(f(\tilde{\mathbf{x}})-f^{\star}\big)
+L\eta^{2}\sigma^{2}.
\tag{18}
\]

\noindent\textbf{[Step 5] Telescoping over the $m$ inner iterations.}
Sum (18) for $t=0,\dots,m-1$ and use $A_{0}=f(\tilde{\mathbf{x}})-f^{\star}$:
\[
\begin{aligned}
&(1-2L^{2}\eta^{2})\sum_{t=0}^{m-1}A_{t+1}
\le \sum_{t=0}^{m-1}A_{t}
-\eta\sum_{t=0}^{m-1}\EE\!\big[\norm{\grad f(\mathbf{x}_{t})}^{2}\big]\\
&\qquad
+2L^{2}\eta^{2}m\big(f(\tilde{\mathbf{x}})-f^{\star}\big)
+L\eta^{2}m\sigma^{2}.
\end{aligned}
\tag{19}
\]
Observe that $\sum_{t=0}^{m-1}A_{t}
= A_{0}+\sum_{t=1}^{m-1}A_{t}
\le A_{0}+ \sum_{t=0}^{m-1}A_{t+1}$.
Hence,
\[
(1-2L^{2}\eta^{2})\sum_{t=0}^{m-1}A_{t+1}
\le A_{0}+\sum_{t=0}^{m-1}A_{t+1}
-\eta\sum_{t=0}^{m-1}\EE\!\big[\norm{\grad f(\mathbf{x}_{t})}^{2}\big]
+\cdots
\tag{20}
\]
Cancel $\sum_{t=0}^{m-1}A_{t+1}$ from both sides, obtaining
\[
-\;2L^{2}\eta^{2}\sum_{t=0}^{m-1}A_{t+1}
\le A_{0}
-\eta\sum_{t=0}^{m-1}\EE\!\big[\norm{\grad f(\mathbf{x}_{t})}^{2}\big]
+2L^{2}\eta^{2}m\big(f(\tilde{\mathbf{x}})-f^{\star}\big)
+L\eta^{2}m\sigma^{2}.
\tag{21}
\]
Drop the non‑positive term $-2L^{2}\eta^{2}\sum A_{t+1}$ to get an inequality in the
desired direction:
\[
\eta\sum_{t=0}^{m-1}\EE\!\big[\norm{\grad f(\mathbf{x}_{t})}^{2}\big]
\le A_{0}
+2L^{2}\eta^{2}m\big(f(\tilde{\mathbf{x}})-f^{\star}\big)
+L\eta^{2}m\sigma^{2}.
\tag{22}
\]

\noindent\textbf{[Step 6] Replace $A_{0}=f(\tilde{\mathbf{x}})-f^{\star}$.}
Thus,
\[
\sum_{t=0}^{m-1}\EE\!\big[\norm{\grad f(\mathbf{x}_{t})}^{2}\big]
\le
\frac{f(\tilde{\mathbf{x}})-f^{\star}}{\eta}
+2L^{2}\eta m\big(f(\tilde{\mathbf{x}})-f^{\star}\big)
+L\eta m\sigma^{2}.
\tag{23}
\]

\noindent\textbf{[Step 7] Choose $\eta\le\frac{1}{3L}$ to bound the second term.}
Since $2L^{2}\eta\le\frac{2}{3}L$, we have
\[
2L^{2}\eta m\big(f(\tilde{\mathbf{x}})-f^{\star}\big)
\le \frac{2}{3}L m\big(f(\tilde{\mathbf{x}})-f^{\star}\big).
\]
Combine with the first term, factor $ \frac{2}{\eta}$:
\[
\sum_{t=0}^{m-1}\EE\!\big[\norm{\grad f(\mathbf{x}_{t})}^{2}\big]
\le
\frac{2\big(f(\tilde{\mathbf{x}})-f^{\star}\big)}{\eta}
+L\eta m\sigma^{2}.
\tag{24}
\]

\noindent\textbf{[Step 8] Average over epochs.}
Let $S$ be the number of completed epochs.
Summing (24) over $s=0,\dots,S-1$ and dividing by $S\,m$ yields
\[
\frac{1}{S\,m}\sum_{s=0}^{S-1}\sum_{t=0}^{m-1}
\EE\!\big[\norm{\grad f(\mathbf{x}_{t}^{\,s})}^{2}\big]
\le
\frac{2\big(f(\tilde{\mathbf{x}}^{\,0})-f^{\star}\big)}{\eta\,S\,m}
\;+\;
\frac{L\eta\sigma^{2}}{S}\sum_{s=0}^{S-1}1.
\]
Since the sum of $1$ over $S$ epochs equals $S$, we obtain exactly bound (5) of
Theorem~\ref{thm:main}.

\noindent\textbf{[Step 9] Optimizing $\eta$.}
Set $\eta=\frac{c}{L\sqrt{S}}$ with a universal constant $c\le\frac{1}{3}$
to satisfy the stepsize restriction.
Plugging into (5) gives
\[
\EE\!\big[\norm{\grad f(\mathbf{x}_{\mathrm{out}})}^{2}\big]
\;\le\;
\underbrace{\frac{2L\big(f(\tilde{\mathbf{x}}^{\,0})-f^{\star}\big)}{c\sqrt{S}}}_{O(1/\sqrt{S})}
\;+\;
\underbrace{4c\,\frac{\sigma^{2}}{\sqrt{S}}}_{O(1/\sqrt{S})}.
\]
Since the total stochastic gradient evaluations equal $N=S(m+n)=\Theta(Sn)$,
the bound can be rewritten as $O(1/N)$ up to the variance term, completing the
proof. \hfill $\blacksquare$

\section*{Rate Derivation (With Constants)}
From (5) and the choice $\eta=\frac{1}{3L\sqrt{S}}$,
\[
\frac{1}{S\,m}\sum_{s,t}\EE\!\big[\norm{\grad f(\mathbf{x}_{t}^{\,s})}^{2}\big]
\le
\frac{6L\big(f(\tilde{\mathbf{x}}^{\,0})-f^{\star}\big)}{\sqrt{S}}
\;+\;
\frac{4\sigma^{2}}{3\sqrt{S}}.
\]
If $m=n$ (a common practical choice), the total number of oracle calls after $S$
epochs is $N=S(2n)$, i.e. $S=\frac{N}{2n}$.  Substituting,
\[
\EE\!\big[\norm{\grad f(\mathbf{x}_{\mathrm{out}})}^{2}\big]
\;\le\;
\frac{12L n}{\sqrt{N}}\big(f(\tilde{\mathbf{x}}^{\,0})-f^{\star}\big)
\;+\;
\frac{4\sigma^{2}\sqrt{2n}}{\sqrt{N}}.
\]
Thus the expected squared gradient norm decays as $O(1/\sqrt{N})$ in the
worst case, but the \emph{deterministic} part (first term) enjoys the faster
$O(1/N)$ behaviour when $\sigma=0$.

\section*{Special Cases}
\begin{itemize}[leftmargin=*]
\item \textbf{Exact finite‑sum ($\sigma=0$).}
      The bound reduces to
      $\displaystyle \frac{2\big(f(\tilde{\mathbf{x}}^{\,0})-f^{\star}\big)}{\eta S m}$,
      i.e. $O(1/N)$.
\item \textbf{Convex $L$‑smooth case.}
      The same derivation yields
      $\displaystyle f(\tilde{\mathbf{x}}^{\,S})-f^{\star}
        \le \frac{2\big(f(\tilde{\mathbf{x}}^{\,0})-f^{\star}\big)}{\eta S m}
        +2L\eta\sigma^{2}$,
      reproducing the classical linear‑convergence for strongly convex
      objectives after adding a strong‑convexity assumption.
\item \textbf{Large inner loop $m\gg n$.}
      When $m$ grows, the variance term $L\eta m\sigma^{2}$ dominates,
      suggesting a practical rule $m=O(n)$ to balance computation and variance.
\end{itemize}

\end{document}